\ChapterImagePrelim[cap:resumen]{Resumen}{./images/fondo.png}\label{cap:resumen}
\mbox{}\\
\noindent
El presente trabajo desarrolla la definición de aspectos de una arquitectura de seguridad para la infraestructura de computación en la nube 
del grupo de investigación en Redes, Información y Distribución (\GRID) de la Universidad del Quindío. La propuesta surge ante la necesidad 
de fortalecer las capacidades de protección, monitoreo y gestión de eventos de seguridad dentro de entornos académicos basados en 
infraestructuras físicas y virtualizadas, considerando lineamientos asociados a estándares internacionales de seguridad de la información y 
computación en la nube.

Metodológicamente, el trabajo comprende: (\textbf{a}) la caracterización contextual de las necesidades, problemas y oportunidades (\NPO) 
del grupo \GRID, (\textbf{b}) el desarrollo de un Estudio de Mapeo Sistemático de Literatura (\SMS) orientado a identificar metodologías, 
arquitecturas y tecnologías relacionadas con seguridad en computación en la nube, (\textbf{c}) la aplicación del proceso de Análisis y 
Resolución de Decisiones (\DAR) del modelo \CMMI~para la selección de tecnologías de seguridad y el diseño e implementación 
de un prototipo funcional basado en un enfoque de defensa en profundidad, (\textbf{d}) la implementación del prototipo funcional y 
(\textbf{e}) la validación del prototipo funcional.

Como resultado, se seleccionaron tecnologías orientadas a Sistemas de Detección y Prevención de Intrusiones (\IDPS) y plataformas de 
Gestión de Información y Eventos de Seguridad (\SIEM), integradas dentro de una arquitectura interoperable y adaptable a contextos 
académicos con recursos limitados. Finalmente, se implementó un prototipo funcional con capacidades de monitoreo continuo, correlación de 
eventos, centralización de registros y respuesta activa ante incidentes, validando la viabilidad técnica y funcional de la propuesta para 
el fortalecimiento de la protección de los activos de información del grupo \GRID.


\ChapterImagePrelim[cap:abstract]{Abstract}{./images/fondo.png}\label{cap:abstract}
\mbox{}\\
\noindent
This work develops the definition of security architecture aspects for the cloud computing infrastructure of the Research Group on Networks, 
Information and Distribution (\GRID) at the Universidad del Quindío. The proposal arises from the need to strengthen protection, monitoring, 
and security event management capabilities within academic environments based on physical and virtualized infrastructures, considering 
guidelines associated with international information security and cloud computing standards.

Methodologically, the work comprises: (\textbf{a}) the contextual characterization of the needs, problems, and opportunities (\NPO) of the 
\GRID\ group, (\textbf{b}) the development of a Systematic Mapping Study (\SMS) aimed at identifying methodologies, architectures, and 
technologies related to cloud computing security, (\textbf{c}) the application of the Decision Analysis and Resolution (\DAR) process from 
the \CMMI\ model for the selection of security technologies and the design and implementation of a functional prototype based on a defense-
in-depth approach, (\textbf{d}) the implementation of the functional prototype, and (\textbf{e}) the validation of the functional prototype.

As a result, technologies oriented toward Intrusion Detection and Prevention Systems (\IDPS) and Security Information and Event Management 
(\SIEM) platforms were selected and integrated into an interoperable architecture adaptable to academic environments with limited resources. 
Finally, a functional prototype was implemented with capabilities for continuous monitoring, event correlation, log centralization, and 
active incident response, validating the technical and functional feasibility of the proposal for strengthening the protection of the 
information assets of the \GRID\ group.
